% <!-- This makes it look pretty in a browser --><pre>
\documentclass[a4paper,english]{cv}
\usepackage[latin1]{inputenc}
\usepackage{a4wide}
\pagestyle{empty}
\usepackage{graphicx}
\usepackage{longtable}
\makeatletter
\newcommand{\HS}{\hspace{0.5cm}}
\date{}
\newcommand{\B}[1]{%
\makebox[1cm][l]{\small{#1:}}}
\newcommand{\VT}{\vspace{2mm}}
\renewcommand{\topicmargin}{7em}
\usepackage{babel}
\makeatother
\begin{document}
\title{Curriculum Vitae}
\maketitle
\thispagestyle{empty}\centerline{\fbox{\begin{minipage}[c]{0.50\columnwidth}%
\begin{center}\hfill{}\hfill{}\begin{tabular}{p{2.1cm}p{4cm}}
Name:&
Martin Strauss\\
Address:&
\parbox[t]{4cm}{Kornhausstrasse 37
\\
CH-8037 Z�rich
\\
SWITZERLAND\VT{}}\\
Telephone:&
\parbox[t]{4cm}{Home: +41 43 500 6897
\\
Mobile: +41 76 260 4905\VT{}}\\
E-mail:& martin@ockle.org\\
Birth:& 19th July, 1983; K�ln, Germany\\
Citizenship:& Australian (naturalised 26th January 1988)\\
\end{tabular}\hfill{}\end{center}\end{minipage}%
\begin{minipage}[c]{0.50\columnwidth}%
~\hfill{}\includegraphics[%
  width=0.50\columnwidth,
  keepaspectratio]{Strauss_Martin.jpg}\hfill{}~\end{minipage}%
}}
\section{Professional Interests}

Software quality and usability: in particular, post-desktop interaction styles such as haptics, non-speech audio, and embodied agents.\section{Skills}
\subsection{Computer Skills}

\begin{topic}
\item [ OS]  Expert user of Unix (especially Linux and SunOS) and Mac OS X, advanced user of Windows.

\item [ Languages]  Extensive experience with C, Java, C++, Prolog, Javascript, PHP, Python and Perl, as well as BASH shell scripting. Also experience with languages such as Haskell, Lex/Yacc, BASIC, RSI IDL and Assembler (primarily x86) as well as some exposure to Pascal, Scheme, Lisp, etc.

\item [ Software Engineering]  Extensive experience with UML and Object-Oriented design, agile development methods, and also with functional and logic programming.  Also experience with Microsoft Visual Studio, Rational Rose, Eclipse, NetBeans.

\item [ Administration]  Extensive experience administering Linux servers (primarily Debian GNU/Linux and Gentoo, but also Red Hat, Mandrake and Slackware), and some experience administering Windows NT4, 2000 and XP workstations and Windows NT4, 2000 and 2003 servers. Experience with directory services (LDAP), SQL databases (MySQL and PostgreSQL), firewalls and routing (netfilter/iptables and ipfilter), Internet proxy services (squid, dante-socks), mail services (sendmail, exim, postfix), etc.

\item [ Tools]  Experience using and programming MATLAB and RSI IDL, and toolboxes such as Simulink.

\end{topic}
\subsection{Language Skils}

\begin{topic}
\item [ English]  native language

\item [ German]  native language

\item [ Japanese]  some knowledge

\end{topic}
\section{Employment}

\begin{topic}
\item [2007 --] \textbf{Software Engineer, Google Inc.}  Freigutstrasse 12, CH-8002 Z�rich, Switzerland

As an engineer on the Google Maps team, I develop the next-generation technologies that change how millions interact, making the world's information about places in the real world accessible and useful to users. I use ideas from every area of computer science, including information retrieval, artificial intelligence, machine learning, natural language processing, distributed computing, large-scale system design, networking, security, data compression and user interface design to handle information at the scale of the interface between the web and the real world.

\item [2006 -- 2007] \textbf{Research Assistant, Deutsche Forschungszentrum f�r K�nstliche Intelligenz}  Stuhlsatzenhausweg 3 (Geb�ude D3 2)
D-66123 Saarbr�cken, Germany

I worked as a research assistant (wissenschaftliche Hilfskraft) at the German research center for AI (DFKI) in the VirtualHuman group, developing software tools for creating, deploying and evaluating the embodied agents ("Virtual Human") developed by the group.

\item [2006] \textbf{Software Engineer, Defence Science and Technology Organisation}  506 Lorimer St, Fishermans Bend VIC 3207, Australia

From January until March I worked as a graduate software engineer in the Air Operations Simulations Center in the Defence Science and Technology Organisation.  The Air Operations Simulations Center operates a facility in which cutting-edge technology is used to produce a realistic flying environment for human-in-the-loop simulation experiments.  The facility conducts research into aerospace operations requirements and equipment suitability for Australian Defence aerospace operations.

\item [2005] \textbf{Tutor, Ormond College}  College Crescent, Parkville, VIC 3052, Australia

During the second semester of 2005, I tutored a group of first year students at Ormond College who were taking the subject `Electrical Circuits'.  In their end-of-term feedback, the students wrote that I `helped [them] to raise questions, express uncertainties and participate fully', and that the tutorials were `stimulating and thought provoking'.

\item [2004 -- 2005] \textbf{Software Engineering Project and Advanced Software Engineering Project}  

As part of my Software Engineering degree at the University of Melbourne, I participated in two industrial projects.  In the first, I was in a team of 5 software engineers who developed a software package for researchers in the Optics group in the School of Physics at Melbourne University.  For the second, I was in a team of 15 software engineers who developed a software package for researchers in the Department of Information Systems at Melbourne University.  In each project, We conducted all of the activities of the software and project lifecycle, including planning and managing the software process, requirements elicitation, documentation, requirements analysis, specification, design, implementation and testing.  As part of both projects, we were required to communicate effectively both with each other and with the client, and undertake a variety of leadership positions during the course of the project.  As technical manager and Executive Officer of the first team, my tasks included training the other team members in the tools and software used during the project, supervising team ``work sessions'', and setting and monitoring weekly team goals and objectives; as quality assurance manager and build manager in the second team, my tasks included setting quality objectives for our product, developing, documenting, monitoring and improving quality  processes for our team, and overseeing and managing the tasks and schedule for the final build of our software system.

\item [2001 -- 2005] \textbf{IT Administrator, Ormond College}  College Crescent, Parkville, VIC 3052, Australia

I worked in the IT administration department at Ormond College, first as Assistant IT Administrator in 2001, then as IT Administrator from 2002 onwards. As well as administering a network consisting of several Linux servers, a computer room of 15 Windows NT/2000 workstations, 15 employee Windows workstations and a variable number of student workstations (Windows, Mac and Linux), I participated in the planning and implementation of an LDAP-based authentication and mail routing system, and later in the planning and implementation of a full network hardware upgrade from 10Mbps unswitched to 1Gbps backbone and 100Mbps to the desktop, fully switched, including the extension of the LAN to include a number of terrace houses in the surrounding suburbs owned by Ormond College.

\item [Summer 1999, 2000] \textbf{Laboratory Assistant, Rohm and Haas Australia Pty. Ltd.}  Hays Road, Point Henry, Geelong, VIC 3220, Australia

During the 1999/2000 and 2000/2001 summer holidays, I worked for Rohm and Haas Australia Pty Ltd in their Geelong Technical Centre. I assisted industrial chemists with research and technical service projects, working at the cutting edge of paint technology.

\end{topic}

\section{Other positions held}

\begin{topic}
\item [2003] \textbf{Head of Pleasant Sunday/Wednesday Evening Subcommittee of the Ormond College Students' Club}

As head of the PS/WE subcommittee, I was responsible for organising a series of relaxed social events throughout the academic year for the members of the Ormond College Students' Club.  I was required to write and work with a budget for these events, and manage a subcommittee of around 20 members to organise venues, food and beverages, and entertainment.

\item [2003 -- 2005] \textbf{Senior Chorister, the Choir of Ormond College}

In addition to singing with the Choir of Ormond College, as Senior Chorister I
had the responsibility of assisting the Director of the Choir in a variety of ways, including managing the Friends of the Ormond Choir and writing an annual newsletter; organising choir events, a choir contact list, etc; assisting the Director in organising concerts and a biennial overseas concert tour; organising the Choir's rehearsals and occasionally leading the Choir in rehearsals; and conducting the Choir in the Director's absence.

\item [2000] \textbf{School Music captain, Geelong Grammar School}

As co-School Music Captain of GGS, I was responsible for assisting the Director of Music in the running of the Music School and musical events, including leadership on the 2000 Choir concert tour of Europe; leadership and organisation of musical groups; MC-ing at various events and functions; and organisation and occasional leadership of rehearsals.

\end{topic}

\section{Professional Associations}

\begin{topic}
\item [1995 -- 2001] \textbf{Member of the CSIRO Double Helix Science Club} 

\item [2001 -- 2005] \textbf{Member of Melbourne University Young Engineers} 

\item [2001 --] \textbf{Student Member of the Association of Professional Engineers, Scientists and Managers} 

\item [2001 --] \textbf{Student Member of the Institute of Engineers, Australia} 

\end{topic}
\section{Education}

\begin{topic}
\item [2006 -- 2007]\textbf{Master of Science} The University of Saarland and Max Planck Institute for Computer Science

Awarded International Max Planck Research School Scholarship.
Masters thesis: ``Realtime generation of multimodal affective sports commentary for embodied agents''.
Graduated with honours.

\item [2004 -- 2005]\textbf{Taste of Research Summer Scholarship} The University of New South Wales and National ICT Australia

I investigated the technologies and applications of the field of affective computing (the use of emotions in AI and interface design).  I produced a preliminary ontology of affective computing concepts, for use in a hypothetical agent or affective application.  The results of my work were presented at a poster session at UNSW, and in a report titled ``The Logics of Emotion: a survey of the field of Affective Computing''(see Publications below, \emph{{[}2{]}})

\item [2001 -- 2005]\textbf{Bachelor of Engineering (Software), Diploma of
    Music (Practical)} The University of Melbourne

Graduated 2005 with first class Honours.

\item [1999 -- 2000]\textbf{Melbourne University Program for High Achieving Students} The University of Melbourne

First year Music Performance, Piano, 1999 (first class Honours); and First Year Mathematics, 2000 (second class Honours, division A)

\item [1999 -- 2000]\textbf{International Baccalaureate} Geelong Grammar School

Score of 44 out of possible 45 with perfect subject scores in 5 of 6 subjects, as well as perfect scores for the Extended Essay and Theory of Knowledge.

Title of Extended Essay: ``An investigation into matrix encryption for the purpose of enciphering text messages''(see Publications below, \emph{{[}1{]}})

\item [1994 -- 2000]\textbf{Geelong Grammar School} \item [1989 -- 1993]\textbf{The Geelong College} \end{topic}

\section{Other Training}

\begin{topic}
\item [December 2004] The Australian Logic Summer School

As part of my National ICT Australia Summer Scholarship, I attended the Australian Logic Summer School at the Australian National University's Research School of Information Sciences and Engineering.  There, I participated in lectures, workshops and seminars on the mathematical and philosophical foundations of logic, and a wide variety of applications of logic in computing, including theorem proving, constraint programming, artificial intelligence and knowledge representation and reasoning.

\item [February 2002] SecureCON

I attended the inaugural University of Melbourne Local IT Experts Security Convention, and participated in seminars and workshops on topics such as securing webservers, using intrusion detection systems, using network traffic monitors, writing a information security policy and managing a security compromise.

\item [April 2000] Australian Business Week

I was one of 66 Victorian students awarded a scholarship to participate in the Australian Business Week. This is a one-week intensive residential program aimed at developing the entrepreneurial and small business culture of young Australians. During this program, I also participated in a Team Skills Development workshop.

\end{topic}
\end{document}
%</pre><!-- This makes it look pretty in a browser -->
